\section{Constrained-KP normalization and canonical flow matching}
\label{app:KP-normalization}

We justify the matrix wave normalization used in
\cref{prop:KP-hierarchy} and explain how its times acquire canonical
Hamiltonians on a fixed birational BA level.

\subsection{Wave normalization, commutativity and source evolution}

At each puncture, the ordinary column has leading vector $e_k$ in
\cref{eq:general-BA-asymptotics}. Collecting its $N$ local Laurent
expansions gives a matrix wave function with identity leading matrix.
It determines a formally invertible dressing operator
$\mathscr W=\Id_N+\mathsf W_1\partial^{-1}+\mathsf W_2\partial^{-2}+\cdots$.
Before the gauge change, the $N$ Laurent expansions of $h_A$ form a
time-independent diagonal formal operator $\mathsf H_A(\partial)$.
The spectral equations at all punctures imply
$\mathcal A_A\mathscr W=\mathscr W\mathsf H_A$. Since the diagonal
operators commute, so do $\mathcal A_L$ and $\mathcal A_M$.
This proves \cref{eq:KP-matrix-commutation} without inferring
commutativity from only one eigenfunction column.

Now use the local coordinates $s'_k=1/h_L(P)$. The gauge
$\mathsf T^{-1}=e^{x\mathsf C_L}$ changes the leading ordinary
factor from $e^{x/s_k}e_k$ to $e^{xh_L(P)}e_k$; the remaining
factor has identity leading matrix. In these coordinates the
renormalized dressing operator therefore satisfies
\[
 \mathcal L_{\rm KP}=\mathscr W\partial\mathscr W^{-1}.
\]
For a higher time, the principal parts of $h_L^n$ prescribe the BA
exponentials in the original $s_k$ coordinates. Multiplying component
$k$ by $\exp(x_n\operatorname{reg}_{R_k}h_L^n)$ supplies the
required regular constant. The dual component receives the inverse
factor. Hence the normalization includes both the Abel velocity and
the component phases.

The BA construction with these times gives the Sato equation
\[
 \partial_{x_n}\mathscr W
 =-(\mathcal L_{\rm KP}^{\,n})_-\mathscr W,
\]
where $-$ retains the negative differential powers. Equivalently,
subtracting $\mathcal B_n=(\mathcal L_{\rm KP}^{\,n})_+$ from the
ordinary BA derivative leaves zero leading values at every puncture.
The fixed divisor and BA uniqueness give the first equation in
\cref{eq:KP-source-flows}.

For the dual, the residue pairing uses $d\Omega_0$.
Converting to the local spectral measure $dh_L$ multiplies the dual
wave by the fixed scalar factor $d\Omega_0/dh_L$. This ratio has
leading value one at every puncture. The residue pairing identifies
the resulting dual wave with the formal adjoint wave operator,
giving $-\mathcal B_n^{\rm ad}$ as its time generator. Since the
factor is independent of all BA times, the same differential
equation holds for $\widetilde{\boldsymbol\psi}^{\vee}$ itself.
These are the standard BA and adjoint-wave arguments
\cite{Krichever1977,KricheverShiota2011} in the present normalization.

For a matrix differential operator $B$ and column functions $F,G$,
the finite-rank identity is
\[
 [B,F\partial^{-1}G^{\trans}]_-
 =B(F)\partial^{-1}G^{\trans}
  -F\partial^{-1}(B^{\rm ad}(G))^{\trans}
\]
\cite[equations (7)--(8)]{ChvartatskyiSydorenko2013}.
It proves preservation of the negative parts of both operators in
\cref{eq:KP-operator-flows}. Their nonnegative parts remain fixed:
for $\mathcal L_{\rm KP}$ this follows from
$[\mathcal L_{\rm KP}^{\,n},\mathcal L_{\rm KP}]=0$; for
$\mathcal M_{\rm KP}$ it follows from
$[\mathcal L_{\rm KP}^{\,n},\mathcal M_{\rm KP}]=0$ and its constant
differential part $\Id_N\partial+\mathsf C_M-\mathsf C_L$.
The commutator of a negative-order operator with either differential
part has negative order. This completes the operator and source-flow
proof for arbitrary $N$.

\subsection{Selecting the canonical generator on a fixed level}

Fix a birational prescribed BA level of genus $g$. By
\cref{thm:BA-higher-flows}, select $g$ spectral Hamiltonians
$H_1,\ldots,H_g$ with independent Abel vectors
$\boldsymbol U^{(1)},\ldots,\boldsymbol U^{(g)}$.
For comparison with the physical Hamiltonian convention write
$x_n=-it_n$, and let $\boldsymbol U_{{\rm KP},n}$ be defined by
$z(t_n)=z_0-it_n\boldsymbol U_{{\rm KP},n}$.
The coefficients of the required spectral combination solve
\begin{equation}\label{eq:KP-Abel-linear-match}
 \sum_{\alpha=1}^{g}\beta_{n\alpha}\boldsymbol U^{(\alpha)}
 =\boldsymbol U_{{\rm KP},n}.
\end{equation}
The matrix of these vectors is invertible. When $g=0$, no spectral
combination is required.

The remaining difference of the puncture time representatives has
zero residue pairing with every holomorphic differential. The
Mittag--Leffler criterion therefore makes it the principal part of
a meromorphic function supported at the punctures, up to an
additive constant. Its evaluations at the ordinary points and its
regular puncture constants give fixed vertex phase differences
$\delta\varphi_\ell^L$, $\delta\varphi_m^M$, $\delta\varphi_k^N$.
The BA theta translations already agree; these vertex phases account
for the remaining exponential factors in the mode arrays.

The power charges generate them through
\begin{equation}\label{eq:KP-canonical-generator}
\begin{aligned}
 \widehat H_n={}&\sum_{\alpha=1}^{g}\beta_{n\alpha}H_\alpha
 -\sum_{\ell=1}^{L}\delta\varphi_\ell^L Q_\ell^L
 +\sum_{m=1}^{M}\delta\varphi_m^M Q_m^M
 +\sum_{k=1}^{N}\delta\varphi_k^N Q_k^N.
\end{aligned}
\end{equation}
The resulting phase rates of $u_{\ell m},v_{mk},w_{k\ell}$ are
respectively $\delta\varphi_\ell^L-\delta\varphi_m^M$,
$\delta\varphi_m^M-\delta\varphi_k^N$ and
$\delta\varphi_k^N-\delta\varphi_\ell^L$; partner rates are opposite.
Thus \cref{eq:KP-canonical-generator} matches every mode, including
its phase. All coefficients are constants after choosing the level,
and the Hamiltonians and charges in the combination commute.
Allowing these coefficients to vary with spectral integrals inside
one Hamiltonian would introduce additional derivatives and is not
the construction used here.
